\section{Development Process and Lessons Learned}
\label{sec:process}

The project began with a broader question: can the pressure identity in
\cref{eq:pressure-identity} be turned into a training regularizer or coupling rule for low-NFE
flow matching? Several variants were implemented and tested, including finite-difference
Lagrangian consistency penalties, Jacobian-vector-product material-derivative penalties,
directional regularization, pressure-aware minibatch coupling, and Sinkhorn versions of those
couplings. These studies were useful but did not produce a robust paper-level training method.
Some variants improved selected two-dimensional benchmarks but failed or became ambiguous on
higher-dimensional mixtures and CIFAR-like settings.

The main lesson was a separation of roles. Data-side pressure profiles can explain why a problem
has curved marginal transport, but the object integrated at inference is the learned neural field.
If the learned field smooths away part of the data-side curvature, an oracle pressure mesh can be
misaligned with the actual ODE being solved. SCTW therefore schedules from the model's own
self-curvature profile. This is less ambitious than pressure-aware training, but it is cleaner:
the method directly asks where the learned ODE is hard for Euler to integrate. Adjoint-SCTW extends
that question from local difficulty to terminal relevance: a locally large defect may be damped by
the remaining flow, whereas a smaller defect may be amplified.

This process also shaped the experimental posture. We compare against uniform Euler because it is
the default low-NFE sampler. We compare against early-, late-, and symmetric-heavy hand meshes
because fixed nonuniform schedules are simple, strong, and sometimes better than SCTW on a given
problem. We report integration error to a high-NFE reference when available, because endpoint
quality can be affected by learned-model bias. The goal is not to hide failures; the goal is to
understand when a self-curvature mesh is useful. A descriptive cross-schedule correlation audit
was added for the same reason. It measures, within each fixed checkpoint and NFE, whether lower
integration error is aligned with the endpoint metric; it is a diagnostic, not an inferential
test over independently sampled solvers.
